\documentclass[11pt,a4paper]{article}
\usepackage[T1]{fontenc}
\usepackage[utf8]{inputenc}
\usepackage{lmodern,amsmath,amssymb}
\usepackage[margin=23mm]{geometry}
\usepackage{booktabs,tabularx,array,xcolor,microtype}
\usepackage[colorlinks=true,linkcolor=teal!60!black,urlcolor=teal!60!black,citecolor=teal!60!black]{hyperref}
\usepackage{fancyhdr,titlesec,needspace}
\titlespacing*{\section}{0pt}{2.0ex plus .6ex minus .2ex}{.9ex plus .2ex}
\pagestyle{fancy}\fancyhf{}
\fancyhead[L]{\small Low-rank Grothendieck rings}
\fancyhead[R]{\small Research announcement}
\fancyfoot[C]{\small\thepage}
\setlength{\headheight}{14pt}
\setlength{\parindent}{0pt}
\setlength{\parskip}{3pt}
\setlength{\emergencystretch}{1.5em}
\newcommand{\FPdim}{\operatorname{FPdim}}
\newcommand{\Vect}{\operatorname{Vec}}
\newcommand{\classcommit}{9891a7df86783eb577a801caa10e0b0f3f84c9be}
\newcommand{\censuscommit}{a35155aae3b190112d2679ad14d2097d13c4dda0}
\newcommand{\classrepo}{https://github.com/sebastienpalcoux/grothendieck-ring-classification}
\newcommand{\censusrepo}{https://github.com/sebastienpalcoux/fusion-ring-census}
\newcommand{\atlas}{https://sebastienpalcoux.github.io/grothendieck-ring-classification}
\newcommand{\prooflink}[2]{\href{\classrepo/blob/\classcommit/#1}{#2}}
\newcommand{\constructionsurl}{https://sebastienpalcoux.github.io/grothendieck-ring-classification/constructions.html}
\title{\vspace{-12mm}\textbf{Low-rank Grothendieck rings:}\\[3pt]
an AI-assisted classification project,\\ constructions, and open problems}
\author{S\'ebastien Palcoux\thanks{Contributor and project coordinator. Beijing Institute of Mathematical Sciences and Applications (BIMSA), Beijing, China. The project's AI provenance and verification status are stated explicitly below.}}
\date{26 October 2026}
\hypersetup{pdftitle={Low-rank Grothendieck rings: an AI-assisted classification project, constructions, and open problems},pdfauthor={Sebastien Palcoux},pdfsubject={Revised research announcement; repository snapshot checked on 26 October 2026}}
\begin{document}
\maketitle
\vspace{-5mm}
\begin{abstract}
We present a public AI-generated research project on fusion rings and their
categorifications over $\mathbb C$, with data extending through rank twelve.
The project reports unrestricted Grothendieck-ring classifications in ranks
two, three, and four, a complete multiplicity-free classification through
rank twelve, and classifications or partial classifications at higher
multiplicities. The released census contains $1{,}434{,}042$ rings in ranks
three through twelve; ten categorification questions remain open within
its stated bounds. The resources include research manuscripts, categorical
models, exclusion arguments, verification code, and an interactive atlas.
One categorical model per ring suffices for our classification problem;
classifying all realizing categories up to equivalence is a further task.
We summarize the current scope, compare it with AnyonWiki, and link to the
online constructions catalogue and unresolved cases.
The project's original research content was generated by AI under human
direction; independent human verification remains pending.
\end{abstract}

\section{The classification question and the current scope}
For a fusion category $\mathcal C$ over $\mathbb C$, write
$b_0=1,b_1,\ldots,b_{r-1}$ for its simple-object classes. Its Grothendieck
ring has multiplication and maximum fusion multiplicity
\[
 b_i b_j=\sum_{k=0}^{r-1}N_{ij}^{k}b_k,
 \qquad N_{ij}^{k}\in\mathbb Z_{\geq0},
 \qquad M(R)=\max_{i,j,k}N_{ij}^{k}.
\]
The rank is $r$, and multiplicity-free means $M(R)=1$. We count fusion rings
up to based-ring isomorphism. All duality types are included, and
commutativity is not assumed. The census enumerates fusion rings without
filtering for categorifiability. The companion classification asks which
of these rings are $K_0(\mathcal C)$ for at least one complex fusion category.

The rank-one ring is $K_0(\Vect)$. In rank two the family is
$X^2=1+nX$, $n\geq0$, and only $n=0,1$ are categorifiable~\cite{Ostrik2}.
Both have $M=1$, since the maximum includes the unit coefficients.
The higher-rank scope is summarized in Table~\ref{tab:scope}.

\begin{table}[htbp]
\centering\small
\setlength{\tabcolsep}{8pt}
\renewcommand{\arraystretch}{1.12}
\begin{tabular}{@{}c c r r r r@{}}
\toprule
Rank & Census bound & Census rings & Excluded & Realized & Unresolved\\
\midrule
2 & all $n\geq0$ & symbolic & $n\geq2$ & 2 & 0\\
3 & $M\leq1{,}000$ & $460{,}353$ & $460{,}348$ & 5 & 0\\
4 & $M\leq128$ & $889{,}530$ & $889{,}515$ & 15 & 0\\
5 & $M\leq20$ & $40{,}999$ & $40{,}982$ & 17 & 0\\
6 & $M\leq9$ & $25{,}421$ & $25{,}362$ & 57 & 2\\
7 & $M\leq5$ & $10{,}006$ & $9{,}976$ & 30 & 0\\
8 & $M\leq3$ & $4{,}500$ & $4{,}403$ & 93 & 4\\
9 & $M\leq2$ & $2{,}225$ & $2{,}145$ & 76 & 4\\
10 & $M=1$ & 194 & 136 & 58 & 0\\
11 & $M=1$ & 182 & 155 & 27 & 0\\
12 & $M=1$ & 632 & 480 & 152 & 0\\
\bottomrule
\end{tabular}
\caption{Reported project status at the pinned snapshot checked on 26 October 2026.
Counts concern the displayed census ranges; they are based-ring counts,
not numbers of tensor-equivalence classes. The project reports
\textbf{unrestricted classification in ranks two, three, and four};
their finite census bounds are not hypotheses of those claims.
All project claims remain subject to independent human verification.}
\label{tab:scope}
\end{table}

The released census contains $1{,}434{,}042$ based-isomorphism classes in
ranks three through twelve: $1{,}433{,}502$ recorded general exclusions,
$530$ realizations, and ten unresolved cases. The project reports a complete
Grothendieck-ring classification at multiplicity one through rank twelve.

The new rank-twelve band contains $632$ rings: $480$ general exclusions
and $152$ realizations. Of the realized rings, $151$ have known unitary
models; \texttt{R12\_1\_623} has a pseudo-unitary realization whose
unitarity is unproved. This is an open question about unitarity, while its
general complex categorifiability is decided.

At higher multiplicities, classification is complete within the displayed
ranges wherever the unresolved count is zero. The complete contiguous
ranges at ranks six, eight, and nine end at $M=3,2,1$, respectively.
There remain two open rings at rank six (one at each of $M=4,8$),
four at rank eight ($M=3$), and four at rank nine ($M=2$).

Here ``realized'' means that a categorical model is recorded, ``excluded''
means that a general obstruction is recorded, and ``complete'' means that
every ring in the stated census range has one of these two decisions.
These labels describe the project's reported mathematical conclusions.
Neither a surviving necessary condition nor a failed construction attempt
decides categorifiability.

\section{Mathematical methods and verification}
The project combines structural arguments with exact finite calculations.
Its criteria include arithmetic and cyclotomic restrictions, formal
codegrees, induction to the Drinfeld center, Frobenius--Schur indicators,
fusion subcategories, module categories, and coherence obstructions.
A positive decision requires a categorical construction together with
an identification of its full fusion tensor with the census representative.
An exclusion must apply to every possible complex categorification;
an obstruction restricted to pivotal, spherical, pseudo-unitary, unitary,
or braided categories retains that restriction.

For example, the
\prooflink{manuscripts/rank4/papers/companion.pdf}{unrestricted rank-four manuscript}
reports exactly fifteen Grothendieck rings, automatic sphericality of every
categorification, and the sharp bounds
\[
 \FPdim(\mathcal C)\leq24+12\sqrt3,
 \qquad M(K_0(\mathcal C))\leq6.
\]
Its argument reduces the problem theoretically to
$\FPdim(\mathcal C)<1000$ and $M\leq31$, then applies exact checks to
$1{,}692$ remaining rings. The pivotal classification of
Dong--Palcoux--Plessis~\cite{DPP} is a separate result with its stated
hypothesis. In higher ranks the project treats general complex
categorification, not only unitary models: the rank-six list includes
a reported spherical realization whose ring admits no pseudo-unitary
categorification.

Within the reported classification through rank twelve, the
\prooflink{docs/SIMPLE_MULTIPLICITY_ONE.md}{simple-ring audit} finds
exactly sixteen nontrivial simple multiplicity-free Grothendieck rings:
$\mathbb Z[C_p]$ for primes $p\leq12$, and
$K_0(\operatorname{PSU}(2)_{2r-1})$ for $2\leq r\leq12$.
Here simplicity means the absence of a proper nontrivial fusion subring,
and $\operatorname{PSU}(2)_k$ is the even-label fusion subcategory of
$\operatorname{SU}(2)_k$.
This extends the rank-six pattern in~\cite[Corollary~5.9]{LPR22};
an analogous question in arbitrary rank was posed
in~\cite[Question~5.10]{LPR22}.

Three verification tasks must be distinguished: checking a released
fusion tensor, checking a categorical argument or certificate, and proving
that enumeration has missed no ring. For the recent rank-seven $M=5$,
rank-eight $M=3$, rank-nine $M\leq2$, and rank-ten through rank-twelve
$M=1$ releases, enumeration completeness relies on the supplied programs
and producer completion records; the full searches were not independently
rerun. Exact tensor-validity and based-isomorphism checks have a narrower
scope. The \prooflink{docs/VERIFY.md}{verification guide} and
\prooflink{evidence/rank12_m1_20261026/README.md}{rank-twelve proof archive}
record the executed checks and their mathematical dependencies.

The \prooflink{evidence/correctness_20261025/README.md}{correctness revision}
repairs several model-identification arguments and computational verification
procedures without changing the earlier classifications. Its records
distinguish current proofs, superseded arguments, completed replays, and
external computations not rerun. Automated checks, separate AI reviews,
and source checksums do not replace independent human verification of
categorical arguments or enumeration completeness.

\Needspace{14\baselineskip}
\section{Comparison with AnyonWiki}
Our fusion-ring census strictly extends the rank--multiplicity ranges
reported in AnyonWiki~\cite[Table~1]{AW}, including multiplicity-free
rings through rank twelve and larger ranges at higher multiplicities.
Our census is entirely AI-generated under human direction, with independent
human verification pending; AnyonWiki's mathematical data were produced
using software developed without AI assistance~\cite{AW}.
The categorification objectives also differ: for our Grothendieck-ring
classification, one categorical model per ring suffices, whereas
AnyonWiki provides a census of multiplicity-free pivotal fusion categories
through rank seven up to equivalence, including compatible pivotal and
braided structures~\cite{VS26}.\footnote{We retain the term ``census''
and the pivotal qualification used in~\cite{VS26}, whose authors discuss
correctness, completeness, and equivalence explicitly.}
Our Grothendieck-ring analysis also treats higher multiplicities.

AnyonWiki's ring enumeration also includes higher
multiplicities~\cite{AW,VS23}. Our categorification work uses established
human mathematics and, for some positive cases, explicit AnyonWiki data;
the \prooflink{evidence/anyonwiki/README.md}{source and exact-comparison records}
document those inputs.

\section{Constructions and unresolved cases online}
The \href{\constructionsurl}{constructions catalogue} records specific
realizations not, to our knowledge, explicitly described in the literature,
together with their known ingredients and the project's checks of the
required hypotheses and fusion rules. These proposed realizations await
independent verification; \emph{new} is used only in this limited sense,
without claiming novelty up to tensor equivalence or exoticness.

The \href{\atlas/unresolved.html}{unresolved-cases page} maintains the
current list of open rings, with exact fusion tensors, invariants, proof
references, and the status of general, pivotal, spherical, pseudo-unitary,
unitary, and braided categorification.

The \href{\atlas/\#rings}{model catalogue} and
\href{\atlas/records.html}{case directory} now include rank twelve in the
same responsive interface, with filters for ring properties. The
\href{\atlas/theory.html}{theory and definitions page} explains these
properties and the categorical structures under consideration.
Individual records distinguish obstructions from implications between
structures, and statements about a category itself from statements about
a Galois conjugate. The live catalogues may evolve; the counts in this
announcement refer to the fixed snapshot below.

\section{Resources, provenance, and independent review}
The project has three complementary public entry points:
\begin{center}\small
\begin{tabularx}{\textwidth}{@{}l X@{}}
\toprule
Resource & Contents\\
\midrule
\href{\censusrepo}{Fusion Ring Census} & Enumeration programs, canonical fusion tensors,
count tables, and census manuscripts.\\[3pt]
\href{\classrepo}{Grothendieck Ring Classification} & Classification manuscripts,
models, exclusions, decision records, and exact replay materials.\\[3pt]
\href{\atlas/}{Interactive atlas} & Unified ring exploration, constructions and case
records, unresolved questions, definitions, manuscripts, and source evidence.\\
\bottomrule
\end{tabularx}
\end{center}
The \href{\atlas/\#papers}{manuscript index} identifies the current
mathematical accounts, and the
\href{\atlas/review.html}{research-status page} explains their verification
scope. The two repositories and the atlas are also listed in
references~\cite{Census,Classification,Atlas}.

\textbf{Fixed snapshot.}
This revision describes classification commit
\href{\classrepo/tree/\classcommit}{\texttt{9891a7df8678}} and census
commit \href{\censusrepo/tree/\censuscommit}{\texttt{a35155aae3b1}},
both checked on 26 October 2026. The
\prooflink{data/current_publication_revision.json}{current publication manifest}
and \prooflink{data/classification_summary_all_ranks.csv}{all-rank decision table}
record the current mathematical scope, including rank twelve.
Rank-specific source manifests retain the exact census inputs used for
each argument; these bindings are not replaced by the overall snapshot.

\Needspace{6\baselineskip}
\textbf{AI provenance.}
The project's original mathematical arguments, research exposition,
enumeration and verification code, and computational outputs were generated
by AI under the contributor's direction. This announcement was likewise
prepared with ChatGPT. Cited theorems, known categories, and imported data
retain their external authorship. Independent human verification and
Lean formalization of the project remain pending.

Independent review is invited of the classification reductions, the scope
of the exact certificates, categorical constructions and their coherence,
and comparisons with the existing literature. Resolving a listed open
case requires either a categorical model or an obstruction covering all
possible categorifications.

\begin{thebibliography}{99}\small
\setlength{\itemsep}{2pt}\setlength{\parskip}{0pt}
\bibitem{Ostrik2} V. Ostrik, \emph{Fusion categories of rank 2},
Math. Res. Lett. \textbf{10} (2003), 177--183.
\href{https://arxiv.org/abs/math/0203255}{arXiv:math/0203255}.
\bibitem{DPP} J. Dong, S. Palcoux, and A. Plessis,
\emph{Classifying Grothendieck rings of pivotal fusion categories of rank four},
\href{https://arxiv.org/abs/2609.33272}{arXiv:2609.33272} (2026).
\bibitem{LPR22} Z. Liu, S. Palcoux, and Y. Ren,
\emph{Classification of Grothendieck rings of complex fusion categories
of multiplicity one up to rank six},
Lett. Math. Phys. \textbf{112} (2022), article 54.
\href{https://doi.org/10.1007/s11005-022-01542-1}{doi:10.1007/s11005-022-01542-1}.
\bibitem{AW} C. Delaney, R. Razzouk, J. Slingerland, and G. Vercleyen,
\emph{AnyonWiki.org},
\href{https://arxiv.org/abs/2610.30081}{arXiv:2610.30081} (2026).
\url{https://anyonwiki.org/}.
\bibitem{VS26} G. Vercleyen and J. Slingerland,
\emph{On low-rank multiplicity-free complex fusion categories},
Commun. Math. Phys. \textbf{407} (2026), article 231.
\href{https://doi.org/10.1007/s00220-026-05705-2}{doi:10.1007/s00220-026-05705-2}.
\bibitem{VS23} G. Vercleyen and J. K. Slingerland,
\emph{On low rank fusion rings},
J. Math. Phys. \textbf{64} (2023), no.~9, 091703.
\href{https://doi.org/10.1063/5.0148848}{doi:10.1063/5.0148848}.
\bibitem{Census} S. Palcoux (project coordinator),
\emph{Fusion Ring Census}, AI-generated research repository (2026).
\url{https://github.com/sebastienpalcoux/fusion-ring-census}.
\bibitem{Classification} S. Palcoux (project coordinator),
\emph{Grothendieck Ring Classification}, AI-generated research repository (2026).
\url{https://github.com/sebastienpalcoux/grothendieck-ring-classification}.
\bibitem{Atlas} S. Palcoux (project coordinator),
\emph{Fusion \& Grothendieck Ring Atlas}, public project website (2026).
\url{https://sebastienpalcoux.github.io/grothendieck-ring-classification/}.
\end{thebibliography}
\end{document}
